\documentclass[11pt]{article}

\usepackage[margin=1in]{geometry}
\usepackage{amsmath,amssymb,amsthm,mathtools}
\usepackage[hidelinks]{hyperref}

\newtheorem{theorem}{Theorem}[section]
\newtheorem{lemma}[theorem]{Lemma}
\newtheorem{proposition}[theorem]{Proposition}
\theoremstyle{remark}
\newtheorem{remark}[theorem]{Remark}

\newcommand{\Arg}{\operatorname{Arg}}

\title{The Triangle--$C_{4k+1}$ Residual Family for Bicyclic Cacti}
\author{}
\date{25 July 2026}

\begin{document}
\maketitle

\begin{abstract}
Let $G$ be a connected bicyclic cactus whose cyclic blocks are $C_3$ and
$C_q$, where $q=4k+1\geq5$.  The cycles may meet in one cut vertex, or they
may be vertex-disjoint and joined by a path of arbitrary length; arbitrary
trees may be attached throughout.  If $n=|V(G)|$, we prove the strict bound
\[
 s^+(G)>n+2-\sec\frac{\pi}{q}>n.
\]
The proof eliminates all off-spine trees exactly, separates the shared-vertex
and disjoint Sachs expansions, and compares the resulting phase pointwise
with the phase of the isolated $q$-cycle.  A product-subpartition argument
makes the comparison uniform in the joining path and in every tree message.
This closes the full mixed residual family left by the sharp cactus doubly
nonnegative bound.
\end{abstract}

\section{Statement and setup}

All graphs are finite, simple, and undirected.  For the adjacency eigenvalues
$\lambda_1,\ldots,\lambda_n$ of a graph $G$, put
\[
 s^+(G)=\sum_{\lambda_j>0}\lambda_j^2,
 \qquad
 s^-(G)=\sum_{\lambda_j<0}\lambda_j^2.
\]
The positive and negative square-energy conjecture was proved by Liu, Tang,
and Zhang~\cite{LTZ}.  Akbari, Kumar, Mohar, Pragada, and Zhang proposed the
stronger assertion that $s^+(G)\geq n$ whenever a connected $n$-vertex graph
has at least $n+1$ edges~\cite{AKMPZ}.  The sharp blockwise doubly
nonnegative argument leaves, among bicyclic cacti, a mixed residual pair
$\{3,4k+1\}$.  We settle every realization of that pair.

\begin{theorem}\label{thm:main}
Let $G$ be a connected bicyclic cactus whose only cyclic blocks are $C_3$ and
$C_q$, with $q=4k+1\geq5$.  The two cyclic blocks may have one common cut
vertex, chosen arbitrarily on each cycle, or may be vertex-disjoint and joined
by their unique path, of arbitrary length.  Allow arbitrary trees to be
attached at arbitrary vertices.  If $n=|V(G)|$, then
\begin{equation}\label{eq:main}
 s^+(G)>n+2-\sec\frac{\pi}{q}>n.
\end{equation}
\end{theorem}

Let $H$ be the \emph{spine}: in the shared case it is the union of the two
cycles, and in the disjoint case it is the union of the cycles and their
unique joining path.  Every component off $H$ is a tree attached to exactly
one spine vertex; otherwise the cactus property or the choice of the joining
path would be violated.

For a graph $F$ with positive vertex activities $a=(a_v)$, write
\begin{equation}\label{eq:weighted-Z}
 Z_F(a)=\sum_{M\text{ a matching of }F}
          \prod_{v\text{ unmatched by }M}a_v.
\end{equation}
When all activities equal $t$, abbreviate this to
\[
 Z_F(t)=\sum_M t^{|V(F)|-2|M|}.
\]
The empty graph has partition $1$.  In particular, every partition appearing
below is strictly positive for $t>0$.

\section{Exact removal of attached trees}

Orient each off-spine branch toward its attachment vertex.  If $u\to p$ is
an oriented branch edge and $T_{u\to p}$ is the subtree below $u$, including
$u$, define
\[
 Q_{u\to p}(t)=
 \frac{Z_{T_{u\to p}}(t)}{Z_{T_{u\to p}-u}(t)}.
\]
Splitting a matching according to the status of $u$ gives the exact recursion
\begin{equation}\label{eq:BP}
 Q_{u\to p}(t)=t+\sum_{w\text{ child of }u}
                         \frac{1}{Q_{w\to u}(t)}\geq t>0.
\end{equation}
Indeed, after the child-subtree partitions are extracted, leaving $u$
unmatched contributes $t$, while matching $u$ to a chosen child $w$
contributes $Z_{T_{w\to u}-w}/Z_{T_{w\to u}}=1/Q_{w\to u}$.

Consequently the exact effective activity at a spine vertex $v$ is
\begin{equation}\label{eq:activity}
 a_v(t)=t+\sum_{u\text{ in a branch at }v}
                    \frac{1}{Q_{u\to v}(t)}=t+y_v(t),
 \qquad y_v(t)\geq0.
\end{equation}
All branch eliminations also produce the common positive factor
\begin{equation}\label{eq:K}
 K(t)=\prod_{u\to v,\ v\in V(H)} Z_{T_{u\to v}}(t)>0,
\end{equation}
where the product has one factor for each branch rooted next to the spine.

We record explicitly why this factor remains common after vertices of a
cyclic block are deleted.  If a retained spine vertex $v$ is not incident to
a branch edge used by a matching, that branch contributes its full
$Z_T$; if the edge is used, extraction of the same $Z_T$ leaves the activity
$1/Q$.  This is exactly~\eqref{eq:activity}.  If $v$ itself is deleted by a
Sachs cycle, no branch edge to $v$ remains and the former branch is now a
disconnected full tree, again contributing exactly $Z_T$.  Thus deletion of
one or both cycle vertex sets leaves every factor in~\eqref{eq:K}, with no
hidden quotient.  A tree described by identifying its root with $v$ is
handled by deleting that root and treating the resulting components as these
branches.

From now on every matching partition on a subgraph of $H$ uses the effective
activities~\eqref{eq:activity}.  Set
\begin{equation}\label{eq:ABD}
 A=Z_{H-V(C_3)}(a),
 \qquad
 B=Z_{H-V(C_q)}(a),
\end{equation}
and, only when the cycles are vertex-disjoint, set
\begin{equation}\label{eq:D}
 D=Z_{H-(V(C_3)\cup V(C_q))}(a).
\end{equation}

\section{The two Sachs expansions}

Let $\phi_G(z)=\det(zI-A(G))$ and normalize on the positive imaginary axis by
\begin{equation}\label{eq:Psi}
 \Psi_G(t)=i^{-n}\phi_G(it)=\prod_{j=1}^n(t+i\lambda_j),
 \qquad t>0.
\end{equation}
Grouping the Sachs expansion by its collection $\mathcal C$ of
vertex-disjoint cycle components gives
\begin{equation}\label{eq:Sachs}
 \Psi_G(t)=\sum_{\mathcal C}
  \prod_{C\in\mathcal C}(-2i^{-|C|})
  Z_{G-V(\mathcal C)}(t).
\end{equation}
Here a triangle has multiplier $-2i^{-3}=-2i$, while, because
$q\equiv1\pmod4$, the $q$-cycle has multiplier
$-2i^{-q}=2i$.

After applying the exact tree factorization, the shared-cut-vertex case is
therefore
\begin{equation}\label{eq:shared}
 \boxed{\ \frac{\Psi_G(t)}{K(t)}=Z_H(a)+2i(B-A).\ }
\end{equation}
There is no double-cycle term: the two cycles have a common vertex, so they
cannot both occur in a vertex-disjoint Sachs collection.

In the vertex-disjoint case both cycles may occur, and their multiplier is
\[
 (-2i)(2i)=+4.
\]
Thus in this case, and only in this case,
\begin{equation}\label{eq:disjoint}
 \boxed{\ \frac{\Psi_G(t)}{K(t)}=Z_H(a)+4D+2i(B-A).\ }
\end{equation}
Accordingly define
\begin{equation}\label{eq:R}
 R=\begin{cases}
 Z_H(a),&\text{if the cycles share a cut vertex},\\
 Z_H(a)+4D,&\text{if the cycles are vertex-disjoint}.
 \end{cases}
\end{equation}
In both cases $R>0$, and~\eqref{eq:shared}--\eqref{eq:disjoint} say simply
\begin{equation}\label{eq:normalized-core}
 \frac{\Psi_G(t)}{K(t)}=R+2i(B-A).
\end{equation}
Since the real part is positive, the continuous argument tending to zero at
infinity is
\begin{equation}\label{eq:Theta}
 \Theta_G(t)=\Arg\Psi_G(t)
 =\arctan\frac{2(B-A)}{R}\in(-\pi/2,\pi/2).
\end{equation}

\section{Uniform product-subpartition comparison}

Let
\begin{equation}\label{eq:Zq}
 Z_q(t)=Z_{C_q}(t)
\end{equation}
be the bare signless matching partition of the $q$-cycle.  The following
strict comparison is the core of the proof.

\begin{lemma}\label{lem:product}
For every $t>0$,
\begin{equation}\label{eq:key}
 R>Z_q(t)(B-A).
\end{equation}
Consequently
\begin{equation}\label{eq:phase-comparison}
 \Theta_G(t)<\arctan\frac{2}{Z_q(t)}.
\end{equation}
\end{lemma}

\begin{proof}
Put $S=V(C_q)$.  Partition the matchings counted by $Z_H(a)$ according to
whether they use an edge between $S$ and $V(H)\setminus S$.  The matchings
using no such boundary edge factor exactly as
\[
 Z_{C_q}(a|_S)\,Z_{H-S}(a)=Z_{C_q}(a|_S)B.
\]
This remains true in the shared case: the common cut vertex lies in $S$, and
the two triangle edges from that vertex to the remaining triangle vertices
are precisely boundary edges.  Let $E\geq0$ be the total contribution from
matchings using at least one boundary edge.  Then
\begin{equation}\label{eq:partition}
 Z_H(a)=Z_{C_q}(a|_S)B+E.
\end{equation}

Every activity on $S$ has the form $t+y_v$ with $y_v\geq0$.  Since the
matching partition has nonnegative coefficients in its vertex activities,
\[
 Z_{C_q}(a|_S)\geq Z_q(t).
\]
Write $L=Z_{C_q}(a|_S)-Z_q(t)\geq0$.  Equations
\eqref{eq:R} and~\eqref{eq:partition} now give, with the last term present
only in the disjoint case,
\begin{equation}\label{eq:positive-difference}
 R-Z_q(t)(B-A)=E+LB+Z_q(t)A\ (+4D)>0.
\end{equation}
The inequality is strict because $Z_q(t)A>0$.  This proves~\eqref{eq:key}.

If $B-A\leq0$, then~\eqref{eq:Theta} is nonpositive whereas the right side
of~\eqref{eq:phase-comparison} is positive.  If $B-A>0$, divide
\eqref{eq:key} by the positive quantity $RZ_q(t)$ to obtain
\[
 \frac{2(B-A)}R<\frac2{Z_q(t)},
\]
and use strict monotonicity of the arctangent.  Thus
\eqref{eq:phase-comparison} holds without any assumption on the sign of
$B-A$.
\end{proof}

This argument uses neither a continuant estimate nor a bounded connector: all
joining-path and attached-tree dependence is absorbed into positive matching
partitions.  In particular, symbolic checks for the smallest $C_3$--$C_5$
cores may corroborate the inequality but are not certificates needed by the
proof.

\section{The isolated-cycle phase and Coulson integral}

For the isolated $C_q$, formula~\eqref{eq:Sachs} reads
\begin{equation}\label{eq:cycle-Psi}
 i^{-q}\phi_{C_q}(it)=Z_q(t)+2i.
\end{equation}
Its continuous phase is therefore
\begin{equation}\label{eq:theta-q}
 \theta_q(t)=\arctan\frac2{Z_q(t)}.
\end{equation}
We next evaluate its weighted phase integral exactly.

\begin{lemma}\label{lem:cycle}
If $q=4k+1$, then
\begin{equation}\label{eq:cycle-difference}
 s^+(C_q)-s^-(C_q)=-2\left(\sec\frac\pi q-1\right)
\end{equation}
and
\begin{equation}\label{eq:cycle-integral}
 \int_0^\infty t\theta_q(t)\,dt
 =\frac\pi2\left(\sec\frac\pi q-1\right).
\end{equation}
\end{lemma}

\begin{proof}
The eigenvalues of $C_q$ are $2\cos(2\pi j/q)$ for $0\leq j<q$.  For
$q=4k+1$, the positive ones correspond to $j=0,1,\ldots,k$ and
$j=q-k,\ldots,q-1$.  Hence
\begin{align*}
 s^+(C_q)
 &=4\left(1+2\sum_{j=1}^k\cos^2\frac{2\pi j}{q}\right)\\
 &=4\left(1+k+\sum_{j=1}^k\cos\frac{4\pi j}{q}\right).
\end{align*}
The finite cosine-sum identity
\[
 \sum_{j=1}^k\cos(jx)
 =\frac{\sin(kx/2)\cos((k+1)x/2)}{\sin(x/2)},
\]
with $x=4\pi/q$ and $q=4k+1$, simplifies to
\[
 \sum_{j=1}^k\cos\frac{4\pi j}{q}
 =-\frac12-\frac14\sec\frac\pi q.
\]
Therefore
\[
 s^+(C_q)=q+1-\sec\frac\pi q.
\]
Since $C_q$ has $q$ edges,
$s^+(C_q)+s^-(C_q)=\operatorname{tr}A(C_q)^2=2q$, which proves
\eqref{eq:cycle-difference}.

For any graph, with the continuous argument specified in~\eqref{eq:Psi}, the
signed Coulson identity is
\begin{equation}\label{eq:Coulson}
 s^+(G)-s^-(G)
 =-\frac4\pi\int_0^\infty t\Theta_G(t)\,dt.
\end{equation}
It follows by summing the scalar identity for the factors
$t+i\lambda_j$; the relation $\sum_j\lambda_j=0$ cancels the first-order
term at infinity.  Applying~\eqref{eq:Coulson} to~\eqref{eq:theta-q} and
using~\eqref{eq:cycle-difference} gives~\eqref{eq:cycle-integral}.
\end{proof}

\begin{proof}[Proof of Theorem~\ref{thm:main}]
Lemma~\ref{lem:product} gives the strict pointwise inequality
$\Theta_G(t)<\theta_q(t)$ for every $t>0$.  The phase integrals in
\eqref{eq:Coulson} converge, so integration and
Lemma~\ref{lem:cycle} yield
\begin{equation}\label{eq:strict-difference}
 s^+(G)-s^-(G)
 >-2\left(\sec\frac\pi q-1\right).
\end{equation}
The strictness is genuine: the continuous positive difference
$\theta_q(t)-\Theta_G(t)$ is nonzero at every $t>0$.

A connected bicyclic graph has $|E(G)|=n+1$.  Consequently
\[
 s^+(G)+s^-(G)=\operatorname{tr}A(G)^2=2n+2.
\]
Averaging this identity with~\eqref{eq:strict-difference} gives
\[
 s^+(G)>n+1-\left(\sec\frac\pi q-1\right)
        =n+2-\sec\frac\pi q.
\]
Finally, $q\geq5$ implies
$1<\sec(\pi/q)\leq\sec(\pi/5)<2$, so the last quantity is strictly larger
than $n$.  This proves both inequalities in~\eqref{eq:main}.
\end{proof}

\begin{remark}[Residual family]
The sharp cactus doubly nonnegative estimate and the independent Sachs-phase
argument settle all other mixed congruence cases but leave cyclic block
lengths $\{3,4k+1\}$.  Theorem~\ref{thm:main} closes that entire mixed
residual, including a shared cut vertex, every joining-path length, and all
tree attachments.  Its conclusion is strict and is stronger than the AKMPZ
target $s^+(G)\geq n$.
\end{remark}

\begin{thebibliography}{9}

\bibitem{AKMPZ}
S.~Akbari, H.~Kumar, B.~Mohar, S.~Pragada, and S.~Zhang,
\emph{Refinement of a conjecture on positive square energy of graphs},
arXiv:2506.07264, 2025.

\bibitem{LTZ}
Y.~Liu, Q.~Tang, and S.~Zhang,
\emph{The positive and negative square-energy conjecture},
arXiv:2607.18031, 2026.

\end{thebibliography}

\end{document}
